Supplements such as the short paper (abstract), the formulation of the assigned
task or an acknowledgement, all have in common that they begin with a headline
and contain a short text underneath. It would often be excessive to start an
unnumbered chapter for this purpose. For this reason, the template provides the
\texttt{tucsimplesection} environment. An example of this is shown in
Listing~\ref{lst:final:simplesection}.

\lstinputlisting[language={[LaTeX]{TeX}},style=block,escapechar=\!,float,caption={Simplesection},label={lst:final:simplesection}]{code/final_simplesection.en.tex}

Figure~\ref{fig:final:simplesection} shows the result produced by the
produced by the environment. The spacing between heading and content can be
be varied by the optional spacing parameter. By default
\verb+\bigskipamount+ is set.

\begin{figure}
\centering
\fbox{%
\begin{minipage}{0.95\linewidth}
\begin{tucsimplesection}{headline}
\lipsum[1]
\end{tucsimplesection}
\end{minipage}}
\caption{Example of a simple caption}
\label{fig:final:simplesection}
\end{figure}

It is the author's responsibility to force the start of a new page
(\verb+\cleardoublepage+). The \verb+\vspace+ commands are used for this purpose,
centre the text block and heading vertically. The distance to the bottom
is larger than the one at the top. This is visually more consistent.
